\section{Discussion}
\label{sec:discussion}

% -----------------------------------------------------------------------
\subsection{Interpretation of Hypothesis Findings}
\label{sec:discussion:interpretation}

\paragraph{H1 — Multi-hop advantage.}
At 4-hop depth, Infon+GNN achieves 66.0\% accuracy versus flat retrieval's
42.0\%, a gap of 24 percentage points.  The gap is monotonically increasing
with depth: 6 pp at 2 hops, 15 pp at 3 hops, and 24 pp at 4 hops.  This
pattern is precisely what the NEXT-edge graph structure predicts: flat
retrieval must match the entire multi-document chain to a single query
vector, which becomes increasingly unlikely as chains lengthen.  MCTS
traversal instead decomposes the chain into single-hop expansions, each
grounded by SPLADE sparse retrieval, so the probability of finding any
individual link is approximately constant across depth.  The sheaf GNN
prior accelerates this by guiding beam selection toward structurally sound
partial chains — those whose type signatures and polarity patterns are
consistent with the hypothesis frame — rather than exploring the full
branching factor uniformly.  Infon-symbolic's parallel advantage (23 pp
at 4 hops) confirms that the gain derives from structured MCTS traversal,
not from the GNN scorer alone.

\paragraph{H2 — Honest abstention.}
Infon+GNN assigns mean $m(\Theta) = 0.90$ to NEI claims and only 0.13
and 0.17 to \textsc{supports} and \textsc{refutes} claims respectively, with
Spearman $\rho = 0.95$.  The LLM baseline assigns $m(\Theta) = 0.60$ to NEI
and 0.35--0.36 to the other classes, with $\rho = 0.52$.  The interpretation
is structural: the DS vacuous element $\Theta$ is the \emph{correct}
mechanism for abstention on absent evidence.  When no infon in the store
aligns with the claim anchor, every combination step in the DS rule returns
mass to $m(\Theta)$ rather than distributing it to $m_S$ or $m_R$.  The LLM
baseline, by contrast, assigns non-zero confidence to \textsc{supports} or
\textsc{refutes} even on NEI claims because its pre-training induces a prior
over labels that is independent of whether any retrieved context actually
supports the claim.  The 8$\times$ reduction in false-positive endorsements
(3\% vs.\ 24\% for Infon+GNN vs.\ LLM) is a direct consequence of this
structural difference.

\paragraph{H3 — DS calibration.}
Infon+GNN achieves ECE~$= 0.05$ versus the NLI classifier's ECE~$=
0.12$.  This result is not incidental to the DS mechanism: DS mass is
constructed from explicit, compositional evidence sources — polarity of
individual infons, pairwise alignment scores between infon predicates and
query predicates, and hop-length attenuation — rather than from a neural
softmax trained on a fixed label distribution.  Softmax outputs from
discriminative classifiers are systematically overconfident at high
confidence bins, a well-documented failure mode that temperature scaling only
partially corrects~\cite{guo2017calibration}.  DS combination, by contrast,
shifts mass toward $\Theta$ whenever evidence is weak or ambiguous,
producing a confidence signal whose relationship to true accuracy is
structurally enforced rather than learned.  The reliability diagrams
(Figure~\ref{fig:reliability_diagrams}) show that DS mass tracks the
diagonal closely across all bins, while the NLI classifier shows the
characteristic upward bow at high confidence.

% -----------------------------------------------------------------------
\subsection{Limitations}
\label{sec:discussion:limitations}

\paragraph{Prior encoder-collapse null result.}
During preliminary experiments evaluating an earlier implementation using
synthetic text templates, all evaluation cells produced exactly 0.332
accuracy --- chance level for a three-class problem.  Root-cause analysis
revealed that BERT's entity tokenisation collapsed on templated patterns:
distinct named entities in the template generated identical contextual
embeddings, making every retrieval key indistinguishable.  This invalidated
the synthetic evaluation harness for that earlier system.  Infon avoids this
failure in two ways: SPLADE-tiny uses sparse activations over the full
vocabulary rather than dense contextual embeddings, so distinct tokens
produce distinct sparse codes regardless of surrounding template structure;
and the current evaluation uses real Wikipedia, AVeriTeC, and SciFact text
rather than template-generated synthetic corpora.  We report this null result
in the interest of reproducibility and to warn practitioners against
evaluating dense retrievers on templated text.

\paragraph{Schema coverage failures.}
When the schema bootstrap step fails to discover relevant anchor predicates,
claim coverage drops below 80\%.  This is most acute in highly technical
domains — protein structure, financial derivatives, legal codification —
where SPLADE-tiny's MS-MARCO pre-training vocabulary has low coverage of
domain-specific terminology~\cite{formal2021splade}.  In our evaluation, we
observe this failure mode most frequently on SciFact claims involving
sub-cellular mechanisms.  Mitigation requires either a domain-adapted
SPLADE model or a learned manifest pruner that can identify low-coverage
claims early and route them to a fallback retriever.  We leave both to
future work.

\paragraph{Transductive GNN.}
The sheaf GNN is trained on synthgen hypergraphs and applied transductively
at inference time.  The model is not scalable beyond approximately 5,000
nodes without re-training, and transfer to out-of-domain corpus topologies
is not validated beyond the three benchmark datasets used here.

\paragraph{MCTS computational cost.}
The MCTS branching factor is proportional to the number of NEXT-edges in
the local infon neighbourhood.  At 4-hop depth over a 10,000-infon store,
a single claim takes approximately 5 seconds on a modern CPU, which is
acceptable for batch academic evaluation but may be prohibitive for
interactive applications requiring sub-second latency.

% -----------------------------------------------------------------------
\subsection{Future Work}
\label{sec:discussion:future}

Several directions follow naturally from the current limitations.

\begin{itemize}
  \item \textbf{Cross-cassette GNN training.}  The GNN is currently trained
    exclusively on synthgen (synthetic hypergraphs with known ground truth).
    Training on real corpus pairs — e.g., HoVer chains extracted from
    Wikipedia — would validate transfer and potentially improve recall on
    out-of-distribution topologies.

  \item \textbf{Multimodal infons.}  The cassette format encodes
    propositional infons over text predicates.  Extending the schema to
    image captions (via vision-language embeddings) and tabular data
    (via structured query expansion) would allow Infon to reason over
    mixed-modality corpora.

  \item \textbf{Learned manifest pruning.}  The current manifest pruner
    uses a bounding-box heuristic over the SPLADE anchor vocabulary to skip
    cassettes that cannot contain a query anchor.  A learned filter —
    e.g., a small classifier trained on (cassette vocabulary, query anchor)
    pairs — would generalise to domain-shifted queries without hand-tuned
    thresholds.

  \item \textbf{Schema autodiscovery without spectral clustering.}  The
    current bootstrap uses spectral clustering over SPLADE embeddings to
    discover connective predicates.  Replacing this with LLM-guided
    ontology induction — where a language model proposes candidate
    predicates from a small seed set of claim-document pairs — would
    eliminate the clustering hyperparameters and scale more gracefully to
    new domains.
\end{itemize}
